% ============================================================
% 分片：s04a-branes  对应原文 PDF 第 42--51 页（印刷页 42--51）
% 负责章节：§4 A-Branes And D-Modules + 4.1--4.6（§4 标题起）
% 译者：trans-s04   审校状态：初稿
% ============================================================

% ==== 原文 PDF 第 42 页（印刷页 42）====
% （本页上半部分为 §3.11 的末两段，属 s03b-genus1 分片范围；本分片自 §4 标题起。）

\section{A-膜与 D-模}

\begin{keytakeaway}
\textbf{核心物理图像：}A-膜支撑在复辛流形 $X=T^*Y$ 的复拉格朗日子流形（典型情形是谱曲线 $D$）上；沿 $D$ 的每个分支，Higgs 场不再是 $E\to E\otimes K_Y$ 的全纯映射，而被扭化为 $\bar\partial_E\varphi=1\otimes\zeta$，于是复联络 $\mathcal{A}=A+i\phi$ 不再平坦，而具有中心曲率。扭化参数中有一部分——B-场对应的 $\zeta'$，以及极点处的分歧参数 $\alpha^*$——并不由 A-膜的几何决定。
\textbf{数学工具：}谱曲线与 Hitchin 纤维化、由谱覆盖给出的线丛或扭转无关层、$K^{1/2}$ 与 $\mathrm{Spin}_c$ 结构、作为 $K$-理论扭化来源的 gerbe、扭 D-模与仿射形变 $\widetilde{T}^*Y$ 的上同调类 $\zeta\in H^1(Y,\Omega^{1,cl}(Y))$、残差定理，以及“B-场即中心曲率”的对应。
\textbf{主要结论：}(1) A-膜只能典范地给出扭 D-模而非通常 D-模，因为 $K_Y^{1/2}$ 没有典范选择，典范的只是张量积 $K_Y^{1/2}\otimes E$；(2) $D$ 出现极点时必须改用仿射形变，此时残差之和不再为零，而等于 $\zeta_{\mathbb{C}}$；(3) 由全纯性论证可得，输出的分歧参数是输入参数 $({}^L\alpha,{}^L\beta,{}^L\gamma)$ 的线性函数，且 $\gamma^*$ 只依赖 ${}^L\gamma$、$\sigma^*=f({}^L\sigma)$。
\end{keytakeaway}

正如 §2.3.1 所许诺的，我们将更详细地讨论 $\mathcal{M}_H$ 上的 A-膜与 $\mathcal{M}$ 上相应的扭 D-模之间的关系。主要目标是解释式 (3.49) 的意义，该式描述 Hitchin 纤维限制到 $T^*\mathcal{M}$ 后在其「缺失的」除子附近的行为（在该例中这些除子是点，因为 $\dim\mathcal{M}=1$）。

% [存疑: 此处原文写 $T^*\mathcal{M}$（无 tilde），但按 §4.5--4.6 的论证，Hitchin 纤维实际落在仿射形变 $\widetilde{T}^*\mathcal{M}$ 之中，式 (3.49) 讨论的正是极点附近的行为；故疑为作者沿用旧记号。已按图像照抄，不加 tilde。]

为此我们需要谱曲线理论，它也将是 §5 的主要工具。因此我们先对该理论作一个非常简短的概述。

尽管 [KW] 第 11 节中把 A-膜与扭 $\mathcal{D}$-模联系起来的论证从物理观点看相当清楚，但它们在数学上完全是猜想性的。下面的讨论又增加了一层猜想，因为尽管我们所描述的内容在定性上与从 [KW] 所能预期的相符，但它尚未建立在同样清晰明确的物理论证之上。

关于这些问题的其他进路，见 [NZ, Nad, Mo]。

\subsection{谱曲线}

这一想法对 $GL_n$ 或 $SL_n$ 解释起来最简单。设 $(E,\varphi)$ 是 Higgs 丛，其中 $E\to C$ 是秩 $n$ 的向量丛，$\varphi\in H^0(C,K\otimes\mathrm{ad}(E))$ 是 Higgs 场。与此数据相联系，如下在余切丛 $T^*C$ 中定义一条曲线 $D$。
% ==== 原文 PDF 第 43 页（印刷页 43）====
% [跨页接续: 原书第 42/43 页，自「Letting z denote ...」起的句子跨页；译文取完整句于本锚点之后。]
设 $z$ 表示余切丛纤维上的线性函数，则曲线 $D$（称为谱曲线）由方程
\begin{equation}\tag{4.1}
\det(z-\varphi)=0
\end{equation}
定义。投影 $\pi:T^*C\to C$ 限制为 $n$ 重覆盖 $\pi:D\to C$。

一般地，$D$ 光滑且不可约，其亏格为 $g'=n^2g-n^2+1$，其中 $g$ 是 $C$ 的亏格。对通常的 Higgs 丛，$D\subset T^*C$ 是一条完备射影曲线；但对分歧的 Higgs 丛则不然，因为 $\varphi$ 有极点。在后一情形，通过添加若干不属于 $T^*C$ 的点，把 $D$ 完备化为一条射影曲线。

% [Tutor Note: 推导补足。$g'=n^2g-n^2+1$ 由 Riemann--Hurwitz 公式给出：对 $n$ 重覆盖 $\pi:D\to C$，有 $2g'-2=n(2g-2)+\deg R$，其中 $R$ 是分歧除子。Higgs 谱曲线的分歧除子次数为 $2n(n-1)(g-1)$（对 $n=2$ 即二次微分 $\det\varphi$ 的 $4g-4$ 个零点），代入即得 $2g'-2=2n^2g-2n^2$，亦即 $g'=n^2g-n^2+1$。注意此处 $g'$ 是“一般”谱曲线的亏格，与 §4.5 中带一个分歧点的具体例子不同。]

定义线丛 $\mathcal{L}\to D$ 为映射 $(\pi^*(\varphi)-z):\pi^*(E)\to\pi^*(E)\otimes K$ 的余核。更确切地说，$D$ 光滑时 $\mathcal{L}$ 是线丛，但一般情形下可能只是扭转无关层。

于是，对 $C$ 上的 Higgs 丛 $(E,\varphi)$，我们联系以下数据：谱曲线 $\pi:D\to C$ 以及线丛或扭转无关层 $\mathcal{L}\to D$。由这些数据可以重构原来的 Higgs 丛。$E$ 由 $E=\pi_*(\mathcal{L})$ 恢复；而 Higgs 场 $\varphi:E\to E\otimes K$ 同样可由 $\pi:D\to C$ 恢复为 $\pi^*(z)$，这里乘以 $z$ 被理解为从 $\mathcal{L}$ 到 $\mathcal{L}\otimes\pi^*(T^*C)$ 的映射。这样，Higgs 丛 $(E,\varphi)$ 便可由对 $(D,\mathcal{L})$ 恢复。

在这一描述中，Hitchin 纤维化的底 $\boldsymbol{B}$ 参数化可能的谱曲线 $D$。例如，对 $G=SL_2$，定义 $D$ 的方程约化为 $z^2+\det\varphi=0$，故 $\boldsymbol{B}$ 由二次微分 $\det\varphi=-\Tr\varphi^2/2$ 参数化。Hitchin 纤维化是把对 $(D,\mathcal{L})$ 映到 $D$ 的映射。Hitchin 纤维化的纤维 $\boldsymbol{F}$ 由给定谱曲线 $D$ 时 $\mathcal{L}$ 的可能选择参数化。

在这一构造中出现哪些线丛 $\mathcal{L}\to D$？对 $GL_n$，若允许所有 Higgs 丛 $(E,\varphi)$、对 $\det E$ 不加限制，则所有 $\mathcal{L}\to D$ 都可出现。于是 $\mathcal{L}$ 对应于 $\mathrm{Pic}\,D$ 中的任意点。若希望通过求解 Hitchin 方程把 $(E,\varphi)$ 联系到一个局部系统，则 $\det E$ 的次数必须为零，且必须把 $\mathcal{L}$ 限制到 $\mathrm{Pic}\,D$ 的某个特定连通分支。用 Riemann--Roch 作一个小计算可知，为使 $E=\pi_*(\mathcal{L})$ 的次数为零，$\mathcal{L}$ 的次数必须是
\begin{equation}\tag{4.2}
\deg\mathcal{L}=n(n-1)(g-1).
\end{equation}
（这一计算的另一版本见 §4.9。）

% [Tutor Note: 推导补足。设 $\pi:D\to C$ 为 $n$ 重覆盖，$\deg\pi_*(\mathcal{L})=n$ 个截面的“次数”。对 $\pi_*\mathcal{L}$ 与 $\mathcal{L}$ 分别用 Riemann--Roch，并用 $\chi(\pi_*\mathcal{L})=\chi(\mathcal{L})$ 及 $\chi(\mathcal{L})=\deg\mathcal{L}+1-g'$，可得 $\deg\pi_*(\mathcal{L})=\deg\mathcal{L}-n(n-1)(g-1)$。令其为零即得 (4.2)。这解释了为何 $GL_n$ 的 Hitchin 纤维是一个 Prym 型阿贝尔簇的平移，而非整个 $\mathrm{Pic}\,D$。]

对 $SL_n$，必须进一步限制 $\mathcal{L}$，使 $\det E$ 平凡，而不仅是次数为零。

\subsection{与 A-膜的关系}

关于 $\mathcal{L}$ 次数的条件 (4.2) 可以有益地表述如下。设 $K_C^{1/2}$ 与 $K_D^{1/2}$ 分别表示 $C$ 与 $D$ 的典范丛的平方根。则 $\mathcal{L}$ 具有形式 $\mathcal{L}=\pi^*(K_C^{-1/2})\otimes K_D^{1/2}\otimes\mathcal{N}$，其中 $\mathcal{N}\to D$ 次数为零。一个富有提示性的重述方式是：我们可以用
\begin{equation}\tag{4.3}
K_C^{1/2}\otimes E=\pi_*(K_D^{1/2}\otimes\mathcal{N})
\end{equation}
来定义 $E$。

最后这一描述的优点在于它与 A-膜理论密切相关。设 $X$ 是辛流形，$L\subset X$ 是拉格朗日子流形。我们从带实拉格朗日子流形的实辛流形这一一般情形开始。为定义支撑在 $L$ 上的 A-膜 $\mathcal{B}$，我们大致上选取一个平坦向量丛
% ==== 原文 PDF 第 44 页（印刷页 44）====
% [跨页接续: 原书第 43/44 页，句子「a flat vector bundle $V\to L$」跨页。]
$V\to L$，以及 $L$ 上的自旋丛 $\mathcal{W}$。于是在微分几何中，膜 $\mathcal{B}$ 与向量丛 $\mathcal{W}\otimes V\to L$ 相联系。

现在特殊化到 $X$ 为复辛流形、$L$ 为复维数 $d$ 的复拉格朗日子流形的情形。此时，记 $K_L^{1/2}$ 为 $L$ 的典范丛的平方根，$\Omega^{0,i}$ 为 $L$ 上 $(0,i)$-形式层，则 $L$ 的自旋丛具有形式 $\mathcal{W}=K_L^{1/2}\otimes(\oplus_{i=0}^{d}\Omega^{0,i})$。因此，在代数几何中人们把 $\mathcal{B}$ 与 $K_L^{1/2}\otimes V\to L$ 的截面层相联系；与 $\oplus_{i=0}^{d}\Omega^{0,i}$ 作张量积只不过给出该层的 $\bar\partial$ 消解。于是在式 (4.3) 中，右端涉及的恰是支撑在拉格朗日子流形 $D$ 上的 A-膜的定义数据，即线丛 $K_D^{1/2}\otimes\mathcal{N}$。稍后我们将看到左端在何种意义上也是自然的。

% [Tutor Note: 物理图像。复辛流形上支撑在复拉格朗日子流形 $L$ 的 A-膜，其“Chan--Paton 丛”是 $K_L^{1/2}\otimes(\oplus_i\Omega^{0,i})\otimes V$；其中 $\oplus_i\Omega^{0,i}$ 只提供 Dolbeault 消解，故在代数几何层面只剩 $K_L^{1/2}\otimes V$。这正是 (4.3) 右端 $K_D^{1/2}\otimes\mathcal{N}$ 的含义：$\mathcal{N}$（次数为零）就是 A-膜上的平坦 $U(1)$ 部分。]

\par\medskip\noindent\textit{注 4.1. 更内蕴的表述避免在 $D$ 与 $L$ 上选取自旋丛。在最精确的描述中，$V\to L$ 并非平坦向量丛，而是平坦扭向量丛。扭化由一个 2 阶复 gerbe $\mathcal{G}$ 给出，其局部平凡化对应于 $L$ 在 $X$ 中法丛上的 spin 结构（$\mathcal{G}$ 由 $\mathbb{Z}_2$ gerbe 经嵌入 $\{\pm1\}\subset\mathbb{C}^\times$ 得到）。用 $\mathcal{G}$ 作扭化，与 $X$ 中的膜其实可以用 $X$ 的 $K$-理论来解释这一事实有关。例如见 [FrW]。要点在于：装备了 $\mathcal{G}$-扭向量丛的子流形 $L\subset X$ 定义了 $K(X)$ 中的一个类，而装备了通常向量丛的 $L$ 并不定义这样的类。}

gerbe $\mathcal{G}$ 还将在 §10 中出现。除该节之外，我们不会纳入这一层额外的微妙性。本文大部分篇幅中我们的主要 A-膜例子是支撑在 Hitchin 纤维化的纤维 $\boldsymbol{F}$ 上的膜。在此情形，法丛的行列式平凡，且具有一个也平凡的全局平方根，从而给出 $\mathcal{G}$ 的典范平凡化。

% [Tutor Note: 易错预警。注 4.1 是全文唯一一处提醒“膜其实取值于 $K$-理论”的地方：只有 $\mathcal{G}$-扭向量丛才给出 $K(X)$ 的类。后续各节（除 §10）一律把 $\mathcal{G}$ 平凡化，理由正是 Hitchin 纤维的法丛行列式平凡；若读者在别处见到“法丛行列式不平凡”的情形（如 §10 的 gerbe 讨论），此处给出的平凡化便不再成立。]

被 $\mathcal{G}$ 扭化的秩 1 丛决定了所谓的 $\mathrm{Spin}_c$ 结构。平坦的秩 1 $\mathcal{G}$-扭丛对应于平坦 $\mathrm{Spin}_c$ 结构。在目前的例子中 $L$ 是一条曲线 $D$，形如 $K_D^{1/2}\otimes\mathcal{N}\to D$（其中 $\mathcal{N}$ 次数为零）的线丛自然地对应于平坦 $\mathrm{Spin}_c$ 结构。这就是在不分别引入张量积中任一因子的情况下，内蕴地刻画形如 $K_D^{1/2}\otimes\mathcal{N}$ 的线丛的方式。

\subsection{从 A-膜到扭 D-模的映射}

现在可以解释把 A-膜映到扭 D-模的基本策略。我们先对复辛流形 $X=T^*Y$ 中的 A-膜来做这件事，且设 $Y$ 是一条曲线。

设 $D\subset X$ 是任一光滑曲线，使投影 $\pi:X\to Y$ 限制为有限覆盖 $\pi:D\to Y$。在 §4.1 中，这样的曲线 $D$ 作为 $C$ 上 Higgs 丛 $(E,\varphi)$ 的谱曲线出现。尽管这是引出问题的例子，我们却给底曲线取了新名字 $Y$，因为在实际应用中 $Y$ 其实不是原来的曲线 $C$，而是 $C$ 上稳定抛物丛的模空间 $\mathcal{M}$。

由于维数上的原因，$D$ 关于 $X$ 的全纯辛结构是拉格朗日子流形。我们如下定义支撑在 $D$ 上的秩 1 A-膜：
% ==== 原文 PDF 第 45 页（印刷页 45）====
% [跨页接续: 原书第 44/45 页，「We define a rank 1 A-brane supported on $D$ by picking ...」跨页。]
在 $D$ 上选取一个平坦 $\mathrm{Spin}_c$ 结构，我们并不十分内蕴地把它写成线丛 $K_D^{1/2}\otimes\mathcal{N}\to D$ 的选择，其中 $\mathcal{N}\to D$ 是次数为零的线丛。

于是，大致说来，按照式 (4.3)，我们由 $K_Y^{1/2}\otimes E=\pi_*(K_D^{1/2}\otimes\mathcal{N})$ 定义向量丛 $E\to Y$。并通过 $\varphi=\pi^*(z)$ 定义 Higgs 场 $\varphi$。于是 $(E,\varphi)$ 是 $Y$ 上的 Higgs 丛。

接下来是关键的一步。通过求解 Hitchin 方程，我们把 Higgs 丛 $(E,\varphi)$ 与 $Y$ 上秩 $n$ 的局部系统联系起来，后者我们也记作 $E$。因此，与其只把 $K_Y^{1/2}\otimes E$ 看作向量丛，我们可以把它的截面层看作 $Y$ 上的扭 D-模，即作用在 $K_Y^{1/2}$ 的截面上的微分算子层 $\mathcal{D}^*$ 的模层。

这里一个重要细节是：从由平坦 $\mathrm{Spin}_c$ 丛（我们粗略地写作 $K_D^{1/2}\otimes\mathcal{N}$）定义的 A-膜出发，我们并不能十分内蕴地得到次数为零的向量丛 $E\to Y$；得到的只是次数非零的向量丛，我们把它非内蕴地写作 $K_Y^{1/2}\otimes E$。$K_Y^{1/2}$ 没有典范选择，因此也没有定义 $E$ 的典范方式。唯有二者的张量积是典范的；这正是为什么 $T^*Y$ 中支撑在曲线 $D$ 上的 A-膜典范地不映到 $Y$ 上的通常 D-模（例如局部系统 $E$ 的截面层），而映到扭 D-模（$K_Y^{1/2}\otimes E$ 的截面层）。

% [Tutor Note: 易错预警。$E$ 只在相差一个 2 阶线丛的扭化下有意义。切勿把“扭 D-模”理解成“$\mathcal{D}$-模再乘一个线丛”：$\mathcal{D}^*$ 与 $\mathcal{D}^*_\chi$ 的区别在于算子作用的对象是 $K_Y^{1/2}$ 还是 $K_Y^{1/2}\otimes T_\chi$（见 §4.6）。也正是这一层扭化，使得 (4.3) 的左端虽无典范分解、却是典范的对象。]

尽管 $E$ 只在相差一个 2 阶线丛的扭化下被定义，这一扭化不影响 Hitchin 方程，且由该构造得到的扭 D-模只依赖于原来的 A-膜，而不依赖其他任何选择。真正不够典范的，是用张量积 $K_Y^{1/2}\otimes E$ 来描述这个扭 D-模。

\subsection{极点}

为应用于几何 Langlands，必须把上述内容推广到曲线 $D\subset T^*Y$ 不紧、且投影 $\pi:D\to Y$ 只是一般地 $n$ 重覆盖的情形。如式 (2.6) 那样，我们假设局部地在某点 $u_0\in Y$ 附近，$\pi:D\to Y$ 的某些分支趋于无穷。设 $u$ 表示 $Y$ 在 $u_0$ 附近的局部坐标，并选取 $X=T^*Y$ 上的第二个坐标 $s$，使辛结构为 $du\wedge ds$。在最基本的情形，我们假设这 $n$ 个分支由函数 $s_i(u)$ 描述，其奇点在 $u=u_0$ 处是单极点。于是
\begin{equation}\tag{4.4}
s_i(u)\sim\frac{c_i}{u-u_0},\qquad i=1,\dots,n
\end{equation}
其中 $c_i$ 是复常数，其中一些可能为零。

设 $\mathcal{U}$ 是点 $u_0\in Y$ 的一个小邻域，并在 $\pi^{-1}(\mathcal{U})$ 上平凡化 $\mathcal{L}$。则 $E=\pi_*(\mathcal{L})$ 在 $\mathcal{U}$ 上自然平凡化；相对于这一平凡化，Higgs 场 $\varphi$ 是对角的，$\varphi=\mathrm{diag}(\varphi_1,\dots,\varphi_n)$，其中 $\varphi_i$ 是在 $u=u_0$ 处有极点的微分：
\begin{equation}\tag{4.5}
\varphi_i\sim c_i\frac{du}{u-u_0}.
\end{equation}

正如 $D$ 紧的情形，支撑在 $D$ 上的 A-膜仍给出 $Y$ 上的 Higgs 丛 $(E,\varphi)$（如上按 $K_Y^{1/2}$ 扭化）。然而，现在 $\varphi$ 在 $Y$ 中那些使 $\pi:D\to Y$ 不是 $n$ 重覆盖的点（例如 $u_0$）处有单极点。
% ==== 原文 PDF 第 46 页（印刷页 46）====

我们仍可求解 Hitchin 方程并把这一数据联系到一个（扭）局部系统，因为 Hitchin 方程在有指定极点时表现良好 [Sim1]。在式 (3.4) 中我们已经描述了进入 Hitchin 方程的场的适当极点行为。但现在我们用 $\alpha^*$、$\beta^*$、$\gamma^*$ 表示刻画曲线 $Y$ 上分歧 Higgs 丛的参数：
\begin{equation}\tag{4.6}
\begin{aligned}
A&=\alpha^*\,d\theta+\dots\\
\phi&=\beta^*\frac{dr}{r}-\gamma^*\,d\theta+\dots.
\end{aligned}
\end{equation}
（这里 $u-u_0=re^{i\theta}$。）在秩 $n$ 情形，参数 $\alpha^*$、$\beta^*$、$\gamma^*$ 取值于 $GL_n$ 紧实形式的极大环面的 Lie 代数；因此它们是虚的对角矩阵。\footnote{紧群的实 Lie 代数的元素在酉表示中作为反厄米矩阵作用。}

% [Tutor Note: 易错预警。式 (4.6) 中的场记为 $\phi$（直立体，正体笔画），而式 (4.7) 中的 Higgs 场记为 $\varphi$（花体）；原文在两种字体间并不一致，请注意逐处核对。另：$\alpha^*,\beta^*,\gamma^*$ 取值于紧实形式 $U_n$ 的 Lie 代数，故是反厄米（即“虚的”）对角矩阵；下面 (4.8)--(4.10) 中模长与相位的分离完全依赖于此。]

我们把名字 ${}^L\alpha,{}^L\beta,{}^L\gamma$ 保留给原 ${}^LG$ 局部系统在 $C$ 上的分歧参数，而用 $\alpha,\beta,\gamma$ 表示它们在 A-模型中的对应物。到此为止，我们在记号上没有区分 ${}^LG$ 与 $G$ 的分歧参数。这是因为我们一直在讨论经典几何，而同样的事实对 A-模型与 B-模型都适用。此后，更精确一些会有帮助。

正如 §2.1.1 所解释的，具有如式 (4.6) 那样的奇点的 Hitchin 方程的解，视我们如何看待它，或者确定一个在 $u_0$ 处有奇点的 Higgs 丛，或者确定一个有这种奇点的局部系统。该 Higgs 丛有极点
\begin{equation}\tag{4.7}
\varphi\sim\sigma^*\frac{du}{u-u_0},\qquad \sigma^*=\frac{\beta^*+i\gamma^*}{2}.
\end{equation}
该局部系统具有单值
\begin{equation}\tag{4.8}
M^*=\exp(-2\pi(\alpha^*-i\gamma^*)).
\end{equation}

置 $\sigma^*=\mathrm{diag}(\sigma_1^*,\dots,\sigma_n^*)$ 与 $\sigma_i^*=\frac{1}{2}(\beta_i^*+i\gamma_i^*)$ 是有用的。比较式 (4.5) 与 (4.7)，我们看出
\begin{equation}\tag{4.9}
\sigma_i^*=c_i
\end{equation}
换言之
\begin{equation}\tag{4.10}
\beta_i^*+i\gamma_i^*=2c_i.
\end{equation}

因此，给定一条具有我们所考虑的那类极点的曲线 $D\subset T^*Y$，$\beta^*$ 与 $\gamma^*$ 由几何决定。单值 $M^*$ 的本征值的模为 $\exp(2\pi i\gamma_i^*)$，它们由 $D$ 决定。本征值的辐角依赖于 $\alpha_i^*$。$D$ 的几何不决定 $\alpha_i^*$；对 $\alpha_i^*$ 的任意选择，我们都能求解 Hitchin 方程并得到 $Y$ 上的扭 D-模。

% [Tutor Note: 推导补足。由 (4.8) 有 $M^*=\exp(-2\pi\alpha^*)\exp(2\pi i\gamma^*)$；$\alpha_i^*$ 是虚数，故 $\exp(-2\pi\alpha_i^*)$ 的模为 1，于是 $M^*$ 第 $i$ 个本征值的模为 $\exp(2\pi i\gamma_i^*)$，这确实是正实数，因为 $\gamma_i^*$ 也是虚数。注意原文把“模”写成 $\exp(2\pi i\gamma_i^*)$ 而非 $\exp(-2\pi\,\mathrm{Im}\,\gamma_i^*)$，二者在 $\gamma^*$ 为纯虚时一致，这不是笔误，而是作者选用了紧实形式的反厄米约定。]

\subsection{应用于我们的例子}

通过考察熟悉的例子——亏格 1 的曲线 $C$ 上带 1 个分歧点的 $SL_2$ 几何 Langlands 纲领——我们可以对 $\alpha^*$ 得到一些认识。在这个问题中，我们并不期望能由迄今考虑的数据唯一确定 $\alpha^*$；它同时也应依赖于原 ${}^LG$ 局部系统的底层参数 ${}^L\alpha$，这是完全自然的。在这个问题中目前还没有基于物理的确定 $\alpha^*$ 的方法，但相当多的信息来自非常简单的论证。（对 $\alpha^*$ 的猜想性描述见 §4.7 末尾。）

% ==== 原文 PDF 第 47 页（印刷页 47）====
在此情形，$C$ 上抛物 $SL_2$ 丛的模空间是亏格为零的曲线 $\mathcal{M}$，如我们已描述过的。我们把上述讨论中的曲线 $Y$ 取为 $\mathcal{M}$，并把 $D$ 取为 Hitchin 纤维化的纤维 $\boldsymbol{F}$，视作 $\widetilde{T}^*Y\subset\mathcal{M}_H$ 中的曲线。在考察这个例子时，我们稍稍超前于叙述，因为这里的 $\widetilde{T}^*Y$ 是 $Y$ 的余切丛的仿射形变，而在 §4.4 中我们考虑的是未形变余切丛中的曲线。暂时我们忽略这一区别，它将在 §4.6 中被纳入。

在这个几何 Langlands 纲领的例子中有两个分歧 Higgs 丛。输入是曲线 $C$ 上的分歧 Higgs 丛 $(E,\varphi)$，参数为 ${}^L\alpha,{}^L\beta,{}^L\gamma$。输出是曲线 $Y=\mathcal{M}$ 上的分歧 Higgs 丛，参数为 $\alpha^*,\beta^*,\gamma^*$。我们希望把输出参数确定为输入参数的函数。（输出参数先验地可能在 $\mathcal{M}$ 上使所诱导的局部系统奇异的各点处不同。）

在几何 Langlands 中，人们把输入的 Higgs 丛看作确定 $C$ 上一个单值为 $M=\exp(-2\pi({}^L\alpha-i{}^L\gamma))$ 的局部系统。输出是 $\mathcal{M}$ 上的扭 D-模，由单值为 $M^*=\exp(-2\pi(\alpha^*-i\gamma^*))$ 的局部系统代表。从这个观点看，显然若存在几何 Langlands 对偶，则 $\alpha^*-i\gamma^*$ 全纯地依赖于 ${}^L\alpha-i{}^L\gamma$，且不依赖 ${}^L\beta$。后一断言成立，是因为输入局部系统不依赖 ${}^L\beta$，尽管与之相伴的 Higgs 丛依赖它。于是特别地有
\begin{equation}\tag{4.11}
\frac{\partial\gamma^*}{\partial{}^L\beta}=0.
\end{equation}

先验地，并不清楚 $\alpha^*-i\gamma^*$ 只是 ${}^L\alpha-i{}^L\gamma$ 的函数（而不依赖 $C$ 上输入局部系统的模）；但片刻之后就会清楚，在我们的例子中这对 $\gamma^*$ 成立，从而由全纯性对 $\alpha^*$ 也成立。

另一方面，Hitchin 方程把输入局部系统等同于 Higgs 丛 $(E,\varphi)$，而这里自然的全纯参数是 ${}^L\sigma=({}^L\beta+i{}^L\gamma)/2$。对偶把它映到 $\sigma=(\beta+i\gamma)/2$ 的某个倍数，而 $\mathcal{M}_H(G)$ 全纯地依赖于后者。于是 $\sigma^*=(\beta^*+i\gamma^*)/2$ 全纯地依赖于 ${}^L\sigma$，且不依赖 ${}^L\alpha$；特别地，
\begin{equation}\tag{4.12}
\frac{\partial\gamma^*}{\partial{}^L\alpha}=0.
\end{equation}

把最后两段的断言合在一起，意味着输出三元组 $(\alpha^*,\beta^*,\gamma^*)$ 是输入三元组 $({}^L\alpha,{}^L\beta,{}^L\gamma)$ 的线性函数。论证这一点的一种方式从 Cauchy--Riemann 方程出发，它断言 $\beta^*+i\gamma^*$ 全纯地依赖于
% ==== 原文 PDF 第 48 页（印刷页 48）====
% [跨页接续: 原书第 47/48 页，句子「$\beta^*+i\gamma^*$ is holomorphic in ${}^L\beta+i{}^L\gamma$」跨页。]
${}^L\beta+i{}^L\gamma$：
\begin{equation}\tag{4.13}
\begin{aligned}
\frac{\partial\beta^*}{\partial{}^L\beta}&=\frac{\partial\gamma^*}{\partial{}^L\gamma}\\
\frac{\partial\gamma^*}{\partial{}^L\beta}&=-\frac{\partial\beta^*}{\partial{}^L\gamma}.
\end{aligned}
\end{equation}

利用此式与式 (4.11)，我们看出 $\partial^2\gamma^*/\partial{}^L\gamma^2=\partial^2\beta^*/\partial{}^L\beta\,\partial{}^L\gamma=-\partial^2\gamma^*/\partial{}^L\beta^2=0$。由于按式 (4.11)、(4.12)，$\gamma^*$ 不依赖 ${}^L\alpha$ 与 ${}^L\beta$，故有 $\gamma^*=f({}^L\gamma)$，其中 $f$ 是两个 Lie 代数之间一个可能非齐次的线性函数。于是 $\alpha^*-i\gamma^*$ 关于 ${}^L\alpha-i{}^L\gamma$ 的全纯性以及 $\beta^*+i\gamma^*$ 关于 ${}^L\beta+i{}^L\gamma$ 的全纯性蕴含
\begin{equation}\tag{4.14}
(\alpha^*,\beta^*,\gamma^*)=(f({}^L\alpha)+d_1,\,f({}^L\beta)+d_2,\,f({}^L\gamma))
\end{equation}
其中 $d_1$、$d_2$ 是某些常数，而 $f$ 是同一个线性函数。特别地，
\begin{equation}\tag{4.15}
\sigma^*=f({}^L\sigma)+d_2/2
\end{equation}
必须是 ${}^L\sigma$ 的线性函数。

% [Tutor Note: 推导补足。这里的逻辑常被略过：先由 (4.11)、(4.12) 知 $\gamma^*$ 与 ${}^L\alpha,{}^L\beta$ 无关；再用 (4.13) 得 $\partial^2\gamma^*/\partial({}^L\gamma)^2=\partial^2\beta^*/\partial{}^L\beta\partial{}^L\gamma=-\partial^2\gamma^*/\partial({}^L\beta)^2=0$，故 $\gamma^*$ 是 ${}^L\gamma$ 的线性函数；最后用 Cauchy--Riemann 把 $\alpha^*-i\gamma^*$ 与 $\beta^*+i\gamma^*$ 的全纯性翻译为实部/虚部的配对，得到同一个线性函数 $f$ 加常数。注意 $f$ 允许“非齐次”，即允许仿射线性。]

我们可以把这一线性性预言与式 (3.49) 的结果相比较，在那里我们算出了本例中的 $\sigma^*$。（如 §3.9.1 所述，我们只考虑 improper 分支，因为对这些问题而言两个分支是等价的。）由式 (3.49) 可见，在允许 $\sigma^*$ 的两个本征值互换的意义下，有
\begin{equation}\tag{4.16}
(\sigma_1^*,\sigma_2^*)=(\sigma_0,0).
\end{equation}
这里 $\pm\sigma_0$ 是 ${}^L\sigma$ 的本征值。特别地，线性性成立，而且线性函数 $f$ 是齐次的且 $d_2=0$。我们猜想（没有坚实的基于物理的论证）也有 $d_1=0$；此时 $\varrho^*=\alpha^*-i\gamma^*$ 是 ${}^L\varrho={}^L\alpha-i{}^L\gamma$ 的齐次线性函数。于是写作 ${}^L\varrho=\mathrm{diag}({}^L\varrho_0,-{}^L\varrho_0)$（$\varrho$ 是迹零的，因为它取值于 $\mathfrak{sl}_2$）以及 $\varrho^*=(\varrho_1^*,\varrho_2^*)$，式 (4.16) 蕴含
\begin{equation}\tag{4.17}
(\varrho_1^*,\varrho_2^*)=({}^L\varrho_0,0).
\end{equation}

${}^L\varrho_0$ 与 $-{}^L\varrho_0$ 之间的 Weyl 对称性之所以丢失，是因为在 §3.9.1 中我们把两个可能的映射 $\mathcal{M}_H^0\to\mathcal{M}$ 之一选作起点，以便把 $\mathcal{M}$ 上的扭 D-模联系到 $\mathcal{M}_H^0$ 上的 A-膜。

% [Tutor Note: 易错预警。(4.17) 不是规范不变的陈述，只是在 §3.9.1 选定 $\mathcal{M}_H^0\to\mathcal{M}$ 的某一个映射之后才成立；换另一个映射会把 $\varrho_1^*$ 与 $\varrho_2^*$ 对调，即把 ${}^L\varrho_0$ 换成 $-{}^L\varrho_0$。这正是原文紧接着说“Weyl 对称性已经丢失”的意思。切勿把 (4.17) 当作不依赖分支选择的恒等式。]

现在回到具绕 $p$ 的半单单值的分歧 $SO_3$ 局部系统 $E\to C\backslash p$。若选取 $E$ 的全纯（或拓扑）结构在 $p$ 上的一个延拓，则 $E$ 绕 $p$ 的单值 $M$ 作为 $SL_2$（而非仅仅是 $SO_3$）的元素有意义，且共轭于 $\mathrm{diag}(\lambda,\lambda^{-1})$，其中 $\lambda=\exp(-2\pi{}^L\varrho)$。$\mathcal{M}$ 上局部系统的类似单值 $M^*$ 共轭于
\[
\mathrm{diag}\bigl(\exp(-2\pi\varrho_1^*),\,\exp(-2\pi\varrho_2^*)\bigr).
\]
于是按式 (4.17)，在共轭意义下 $\mathcal{M}$ 上局部系统的单值为
\begin{equation}\tag{4.18}
M^*=\mathrm{diag}(\lambda,1).
\end{equation}

以上所述适用于 $\mathcal{M}_H$ 上支撑在 Hitchin 纤维化的光滑纤维 $\boldsymbol{F}$ 上的一般 A-膜。这样的 A-膜与
% ==== 原文 PDF 第 49 页（印刷页 49）====
% [跨页接续: 原书第 48/49 页，句子「Such an A-brane is associated with a rank 2 twisted local system ...」跨页。]
$\mathcal{M}$ 上在少数几个点处奇异、秩 2 的扭局部系统相联系，且各奇点处的单值恰如上所述。它们在每个奇点处实际上都相同。

对几何内窥，我们想要的是 $\boldsymbol{F}$ 为两个分支 $\boldsymbol{F}_1$ 与 $\boldsymbol{F}_2$ 之并时的类似描述。只支撑在其中一个分支上的分数 A-膜与 $\mathcal{M}$ 中某个 Zariski 开集上秩 1 的扭局部系统相联系。注 3.5 中已说明，每个拉格朗日子流形 $\boldsymbol{F}_1$ 或 $\boldsymbol{F}_2$ 恰在 4 个点 $z^*$ 中的 2 个点处有类似于式 (3.49) 的单极点。因此，与这两个膜中任一个相对应的 $\mathcal{M}$ 上秩 1 局部系统恰有 2 个奇点，在每个奇点处单值为 $\lambda$。

这似乎会导致矛盾。秩 1 局部系统的结构群是阿贝尔群 $GL_1$，故其单值之积必须等于 1，而不是我们在上一段似乎得到的 $\lambda^2$。这个问题的解决牵涉到我们一直忽略的一点：$\mathcal{M}_H^0$ 不是 $\mathcal{M}$ 的余切丛，而是其余切丛的仿射形变。如下文所讨论，这导致局部系统的一种扭化，使得其奇点处单值之积是中心的，但并不简单地等于 1。

% [Tutor Note: 物理图像。$\lambda^2\neq1$ 与 $GL_1$ 的阿贝尔性之间的矛盾，是本文提请注意的关键点：正因为 $\mathcal{M}_H^0$ 是仿射形变而非余切丛，秩 1 局部系统被扭化，绕所有奇点的和乐只落在中心里（即只要求乘积为标量），从而“乘积为 1”被替换为“乘积为中心元素”。这正是 (4.20) 中残差和不再为零的几何根源。]

\subsection{中心扭转}

我们下一个任务是理解 $X=\widetilde{T}^*Y$ 中的 A-膜，其中 $X$ 是 $T^*Y$ 的仿射形变。该形变由类 $\zeta\in H^1(Y,\Omega^{1,cl}(Y))$ 分类。把 $\zeta$ 看作闭 $(1,1)$-形式，积分 $\int_Y\zeta/2\pi i$ 是一个复数，我们称之为 $\zeta_{\mathbb{C}}$。

$\widetilde{T}^*Y$ 中的 A-膜与 $Y$ 上的扭 D-模相联系，其扭化方式比本文迄今讨论的更为有趣。相关的 D-模不是作用在 $K_Y^{1/2}$ 上的微分算子层的模，而是作用在 $K_Y^{1/2}\otimes\mathcal{T}_\chi$ 上的这类算子层的模，其中 $\mathcal{T}_\chi$ 是第一 Chern 类为 $\chi$ 的「线丛」。（$\chi$ 是一个其与 $\zeta$ 的关系有待描述的类。）除非类 $\chi$ 是整的，$\mathcal{T}_\chi$ 并不真正作为线丛存在，但作用在 $K_Y^{1/2}\otimes\mathcal{T}_\chi$ 上的微分算子层 $\mathcal{D}_\chi^*$ 在任何情形下都存在。

在几何 Langlands 的分歧情形，这可以由临界水平的仿射 Kac--Moody 代数的表示 [FG, F2]，或通过分析典范余迷向膜 $\mathcal{A}_{cc}$（如 [GW] §4.4 那样）推出。为了从本文所采取的观点重新得到这一结果，我们取 $X$ 中的曲线 $D$，使投影 $\pi:X\to Y$ 限制为一般地有限覆盖的映射 $\pi:D\to Y$。不过，如 §4.4 那样，我们允许 $D$ 在 $Y$ 中有限多个点上方有「极点」。$D$ 还带有合适的线丛 $\mathcal{L}$。先考察 $\pi:D\to Y$ 一般地为 1 次这一情形是有启发的，此时 $D$ 是 $\pi:X\to Y$ 在 $Y$ 的某个 Zariski 开集上的截面。在此情形，若 $\zeta\neq0$，则必定存在极点。否则 $D$ 将是 $\pi$ 的全局全纯截面，其存在会使类 $\zeta$ 平凡化。

在 $X$ 实际上整体就是余切丛 $T^*Y$ 的情形，曲线 $D$ 对应于 $X$ 上的一个亚纯微分。这个微分就是与 $D$ 相对应的秩 1 Higgs 丛的 Higgs 场 $\varphi$。在 $\varphi$ 有极点的每个点 $u_i\in Y$ 处，可以定义其留数，即式 (4.7) 中称作 $\sigma_i^*$ 的量。按留数
% ==== 原文 PDF 第 50 页（印刷页 50）====
% [跨页接续: 原书第 49/50 页，「According to the residue theorem」跨页。]
定理，留数之和为零：
\begin{equation}\tag{4.19}
\sum_i\sigma_i^*=0.
\end{equation}

当我们用 Hitchin 方程（在这个秩 1 情形它们约化为通常的 Hodge 理论）把对 $(D,\mathcal{L})$ 与 $Y$ 上秩 1 局部系统 $E$ 联系起来时，$u_i$ 处单值的绝对值为 $\exp(-2\pi\,\mathrm{Im}\,\sigma_i^*)$。留数之和为零反映为：对 $GL_1$ 局部系统，这些绝对值之积为 1。（单值的相位由 $\alpha_i^*$ 决定。）

现在考察 $X=\widetilde{T}^*Y$ 实际上是 $T^*Y$ 的仿射形变的情形。此时在每个奇点 $u_i$ 处，我们仍可定义留数 $\sigma_i^*$；确实，仿射形变局部地平凡，尽管没有典范方式使其平凡化，但平凡化的改变不影响留数。仿射形变情形的新特点是留数之和不再为零。代之以有
\begin{equation}\tag{4.20}
\sum_i\sigma_i^*-\zeta_{\mathbb{C}}=0.
\end{equation}
（推导这一结果的一种方式是注意到：我们可以在任意选定的点 $r\in Y$ 之外平凡化 $\zeta$。相对于这一平凡化，$D$ 多出一个奇点，留数为 $-\zeta_{\mathbb{C}}$，而式 (4.20) 不过是所有留数之和为零。）（留数之和为非零常数，对应于单值绝对值之积为不等于 1 的常数，事实上等于 $\exp(2\pi\,\mathrm{Im}\,\zeta_{\mathbb{C}})$。）

这解决了 §4.5 末尾提到的疑难。在那次计算中，我们用了 $Y=\mathcal{M}$ 上的坐标 $z$ 与 $\widetilde{T}^*\mathcal{M}$ 纤维上的坐标 $v$。这些坐标的选择在 $z=\infty$ 之外平凡化了仿射形变。特别地，形变在 $D$ 的极点附近被平凡化，因而没有影响那些极点的留数。在这一平凡化下，$D$ 在 $z=\infty$ 处还有一个极点，自然，包括最后这个极点的所有留数之和为零。

在微分几何中，可以把 $\zeta$ 表示为闭 $(1,1)$-形式\footnote{当然，在曲线上每个 $(1,1)$-形式都是闭的，但这一构造在更高维时确实有意义，如 §4.9 所讨论。}，并把这一情形描述如下。$\pi:\widetilde{T}^*Y\to Y$ 的亚纯截面 $D$ 并不对应于亚纯微分 $\omega$，而是对应于 $Y$ 上的 $(1,0)$-形式，具有如上极点，并满足
\begin{equation}\tag{4.21}
\bar\partial\varphi=\zeta.
\end{equation}
对此式在 $C$ 上积分便回到式 (4.20)；积分左端时使用 $\varphi$ 的极点，并把右端的积分用 $\zeta_{\mathbb{C}}$ 表示。

% [Tutor Note: 推导补足。(4.20) 与 (4.21) 是同一事实的两种表述。对 (4.21) 在 $C$ 上积分：左端 $\int_C\bar\partial\varphi$ 由留数定理化为绕 $D$ 各极点的留数之和（即 $\sum_i\sigma_i^*$），右端 $\int_C\zeta$ 按定义等于 $2\pi i\,\zeta_{\mathbb{C}}$；于是得到 $\sum_i\sigma_i^*=\zeta_{\mathbb{C}}$，即 (4.20)。这解释了为何 $\zeta_{\mathbb{C}}$ 必须是常数而非函数：$\bar\partial\varphi=\zeta$ 的可解性要求 $\zeta$ 的 $(1,1)$ 类为零，而此处恰由 $H^1(Y,\Omega^{1,cl})$ 中的类给出这一“障碍”。]

现在考察其秩 $n$ 情形。换言之，我们假设 $D\subset\widetilde{T}^*Y$ 是一条曲线，使 $\pi:D\to Y$ 一般地是 $n$ 重覆盖。$D$ 还带有线丛 $\mathcal{L}$。我们能否把 $GL_n$ Higgs 丛与对 $(D,\mathcal{L})$ 相联系？可以像通常那样由 $E=\pi_*(\mathcal{L})$ 定义丛 $E$，也可以模仿 Higgs 场的通常定义 $\varphi=\pi^*(z)$。然而现在 $\varphi$ 不是全纯映射
% ==== 原文 PDF 第 51 页（印刷页 51）====
% [跨页接续: 原书第 50/51 页，句子「$\varphi$ is not a holomorphic map $E\to E\otimes K_Y$」跨页。]
$E\to E\otimes K_Y$。更确切地说，$\varphi$ 的每个分支——或更准确地说，$\pi:D\to Y$ 每个分支对 $\varphi$ 的贡献——都恰如式 (4.21) 那样被扭化。于是 $\varphi$ 所满足的方程是
\begin{equation}\tag{4.22}
\bar\partial_E\varphi=1\otimes\zeta,
\end{equation}
其中 $\bar\partial_E$ 是丛 $E$ 的 $\bar\partial$ 算子，而 $1:E\to E$ 是恒等映射。

这类扭化对 Hitchin 方程 (2.1) 有意义，我们可以把它稍作推广为
\begin{equation}\tag{4.23}
\begin{aligned}
F-\phi\wedge\phi&=1\otimes i\zeta'\\
d_A\star\phi&=1\otimes i\,\mathrm{Im}\,\zeta\\
d_A\phi&=1\otimes i\,\mathrm{Re}\,\zeta.
\end{aligned}
\end{equation}

除上一段所见到的由 $\zeta$ 给出的形变（它已被用实微分几何重写）之外，我们在第一个方程中还纳入了一个额外的实\footnote{这些方程右端出现因子 $i$ 的原因是：$\phi$ 与 $F$ 是取值于 $U_n$（$GL_n$ 的紧实形式）的 Lie 代数的形式；我们用反厄米 $n\times n$ 矩阵表示这个 Lie 代数。}闭 $(1,1)$-形式 $\zeta'$。在平凡等价的意义下，这些方程只依赖于 $\zeta'$、$\mathrm{Im}\,\zeta$、$\mathrm{Re}\,\zeta$ 的上同调类。这些类是三个实参数，就像分歧 Higgs 丛的通常参数 $(\alpha,\beta,\gamma)$ 一样。

我们预期，$\widetilde{T}^*Y$ 中的 A-膜应当通过求解这一扭化形式的 Hitchin 方程而映到 $Y$ 上的扭 D-模。$\zeta'$ 的正确选择不由 A-膜的几何决定，有点类似于 §4.4 中 $\alpha^*$ 的取值。这一问题将在 §4.7 中讨论。

在解出扭化方程 (4.23) 后，我们构造复值联络 $\mathcal{A}=A+i\phi$。通常 Hitchin 方程蕴含 $\mathcal{A}$ 平坦。然而在目前情形，$\mathcal{A}$ 并不平坦，而是具有中心曲率 $\mathcal{F}_0=i(\zeta'+i\,\mathrm{Re}\,\zeta)$。具有中心曲率 $\mathcal{F}_0$ 的联络描述的不是 D-模，而是扭 D-模，扭化为 $\mathcal{T}_\chi$，即第一 Chern 类为 $\chi=\mathcal{F}_0/2\pi i$ 的「线丛」。在更内蕴的描述中，扭化其实是按 $K_Y^{1/2}\otimes\mathcal{T}_\chi$ 进行的，其中 $K_Y^{1/2}$ 的角色已在前文描述过。

% [Tutor Note: 物理图像。$\mathcal{A}=A+i\phi$ 的中心曲率 $\mathcal{F}_0=i(\zeta'+i\,\mathrm{Re}\,\zeta)$ 是本节的收束点：通常的 Hitchin 方程给出平坦 $\mathcal{A}$，从而给出 D-模；而仿射形变与 B-场把平坦性破坏为一个中心（即标量）曲率，于是得到扭 D-模，其扭化类 $\chi=\mathcal{F}_0/2\pi i$ 正是 §4.7 要由 B-场算出的量。注意 $\zeta'$ 与 $\mathrm{Re}\,\zeta$ 分别贡献 $\mathcal{F}_0$ 的虚部与实部，切勿把二者混同。]
